\documentclass[noauthor]{ximera}

\title{Week 1 Solutions}
\author{Math 1193 Design Team}

\begin{document}
\begin{abstract}
    Solutions for Week 1.
\end{abstract}
\maketitle

\begin{problem}
    Factor the polynomial \(x^2 - x - 6\). Compare with 
your group and make sense of why each of your methods
works. If you all used the same method, try to find
a different one.

Now solve the equation \(x^2 - x - 6 = 0\).

\begin{explanation}
The AC method has you multiply the coefficient of the \(x^2\) and constant terms.
In this case, \(1 \times (-6) = -6\). Now you need to find two numbers that multiply
to \(-6\) and add to \(-1\), which is the coefficient of the \(x\) term.
Two numbers that work are \(-3\) and \(2\). Then we can write
\begin{align*}
x^2 - x - 6 & = x^2 -3x + 2x - 6 \\
& = (x^2 - 3x) + (2x - 6) \\
& x(x - 3) + 2(x - 3) \\
& = (x - 3)(x + 2)
\end{align*}
Take some time and see if you can understand why each equality is true.

In the first line, since \(-1 = -3 + 2\), \(-x = -3x + 2x\) by combining like terms.

In the second line, we did some grouping with parentheses.

In the third line, we factored an \(x\) out of the first set of parentheses and a 2 out of the second.

In the last line, we did some more complicated factoring. Notice that \((x - 3)\) is a factor of both terms, so 
we can pull it out. We're left with \((x + 2)\).

Our equation then can be rewritten as \((x - 3)(x + 2) = 0\). When two factors multiply to give you \(0\), 
either the first factor must equal \(0\), so \(x - 3 = 0\), or the second must equal \(0\), so \(x + 2 = 0\). 

Solving these two equations yields \(x = 3\) and \(x = -2\) as solutions. To check, you can plug these back into the original equation.
\end{explanation}
\end{problem}

\begin{problem}
Factor the polynomial \(x^2 - 25\). Compare with 
your group and make sense of why each of your methods
works. If you all used the same method, try to find
a different one.

Now solve the equation \(x^2 - 25 = 0\).

\begin{explanation}
The difference (subtraction) of squares method only works for expressions of the form \(a^2 - b^2\), but can be useful to save time. If you have
something that looks like \(a^2 - b^2\), it factors as \((a - b)(a + b)\).

Notice that \(25 = 5^2\), so our polynomial becomes \(x^2 - 5^2\). Because it's a difference of squares, it factors as \((x - 5)(x + 5)\).\

Now like in the previous problem, our equation then can be rewritten as \((x - 5)(x + 5) = 0\). When two factors multiply to give you \(0\), 
either the first factor must equal \(0\), so \(x - 5 = 0\), or the second must equal \(0\), so \(x + 5 = 0\). 

Solving these two equations yields \(x = 5\) and \(x = -5\) as solutions. To check, you can plug these back into the original equation.


\end{explanation}
\end{problem}

\begin{problem}
Solve the inequality \(-2x + 4 \ge 0\). Shade the 
solution set on the number line, then write the 
solution set in interval notation.

\begin{explanation}
When we solve an inequality, we want to find \textbf{all} \(x\)-values that make the inequality true. For example, \(0\) is a solution 
to the inequality (try testing it!).

One solution could be
\begin{align*}
    -2x + 4 & \ge 0 \\
    -2x & \ge -4 \\
    x & \le 2
\end{align*}
Take some time and see if you understand each step.

In the first step, we subtracted \(4\) from both sides of the inequality. 

In the second, we divided by \(-2\). Notice that when we divide or multiply by a negative number on both sides of an inequality,
the inequality changes from \(\ge\) to \(\le\). This also happens with \(\le\) to \(\ge\), \(<\) to \(>\) and \(>\) to \(<\).

The number line looks like this. All numbers to the left of 2 are shaded, since \(x \le 2\). The underline on \(\le\) means that 2 itself is included
in the solution set, so 2 is shaded with a closed dot.
\begin{center}
    
\includegraphics[width=10cm]{number-line.png}
\end{center}

Interval notation is a compact way to write down sets on the number line. Since we've shaded all real numbers to the left of 2 (including 2), we use 
the interval \((-\infty, 2]\). The square bracket indicates we include 2, and the \(-\infty\) on the left indicates that we want everything less than 2 as well.
Remember never to use a square bracket on \(\infty\) or \(-\infty\). 
\end{explanation}
\end{problem}

\begin{problem}
Let's play a game (sort of)!

Consider the equation \(y = 2 + 3x\). Pair up and 
    choose one person to be Player X and the other
    to be Player Y.
    \begin{enumerate}
        \item Player X comes up with 3 \(x\)-values, and 
        Player Y comes up with 3 \(y\)-values. 
        These can be any numbers you want, but try 
        to keep them small so they can be plotted 
        later.
        \item The players exchange their numbers. 
        Player X has to find the \(x\)-values that 
        correspond to the \(y\)-values and put 
        them into coordinate pairs, and Player Y 
        has to find the \(y\)-values that 
        correspond to the \(x\)-values and put 
        them into coordinate pairs. 
        
        Whose job is easier?
        
        \item Write down the 6 
        ordered pairs that satisfy the equation.
        \item Graph your ordered pairs on the \(xy\)
        coordinate plane.
    \end{enumerate}


\begin{explanation}
I would really recommend playing the game, even if you need to play by yourself. 

Here's an example game that Karl and Trae played. Karl was Player X and Trae was Player Y. See if you can play the role of each player. 

Karl chose 1, 2, and 3 as his \(x\)-values. Trae chose \(-17\), 0, and 90 as his \(y\)-values. Trae didn't pay attention to the part about using small numbers.

After receiving Trae's \(y\)-values, Karl found \(-19/3\), \(-2/3\), and \(88/3\) as the corresponding \(x\)-values. 

After receiving Karl's \(x\)-values, Trae found 5, 8, and 11 as the corresponding \(y\)-values. 

Karl and Trae agreed Trae's job was easier.

The ordered pairs were \((1, 5)\), \((2, 8)\), \((3, 11)\), \((-19/3, -17)\), \((-2/3, 0)\), and \((88/3, 90)\).

Remember that when graphing a point like \((1, 5)\), you need to find the point whose \(x\)-coordinate is at 1 and whose \(y\)-coordinate is at 5. Here's
the graph Karl and Trae came up with.

\begin{center}
    \includegraphics[width=10cm]{graphing-game.png}
\end{center}


\end{explanation}
\end{problem}

\begin{problem}
Jaden was given the following problem:

\textit{Let \(f(x) = x^2 + x + 1\). Find \(f(a + 2)\).}

Jaden's work is as follows:

\begin{align*}
    f(a + 2) & = a + 2^2 + a + 2 + 1 \\
    & = a + 4 + a + 3 \\
    & = 2a + 7
\end{align*}

Is Jaden's work correct? Why or why not? Where did he 
go wrong? What did he do correctly?

\begin{explanation}
Jaden's work is not correct. He forgot to put parentheses around \(a + 2\) when substituting into the equation. Here is an example of correct work:
\begin{align*}
    f(a + 2) & = (a + 2)^2 + (a + 2) + 1 \\
    & = (a + 2)(a + 2) + a + 2 + 1 \\
    & = a^2 + 2a + 2a + 4 + a + 3 \\
    & = a^2 + 5a + 7
\end{align*}

Try to make sure you understand each step and ask if you have questions!
\end{explanation}
\end{problem}

\begin{problem}
What are the domain and range of the function 
whose graph is shown below? Compare your answer with
a partner and discuss. 

\begin{center}
    \includegraphics[width=10cm]{domainRangeGraphically.png}
\end{center}

\begin{explanation}
The domain is \((0, \infty)\), and the range is \([-5, \infty)\).

One way to find domain is to imagine sliding a vertical line (or a pencil) from left to right on the graph. Where does the line intersect the graph? Where does it not?
The domain is all \(x\)-values where the line intersects the graph.

For the range, imagine sliding a horizontal line from bottom to top on the graph. The range is all \(y\)-values where the line intersects the graph.

Remember that square brackets never go on \(-\infty\) or \(\infty\). 
\end{explanation}
\end{problem}

\end{document}
